% =========================================================
% Consistency
% =========================================================

\section{Consistency}
\label{sec:consistency}

\begin{remark*}[Toolkit: Consistency]
Two elements of a poset are consistent when they share a common upper
bound --- when they can be jointly extended to some element above both of
them.  In lattice-theoretic settings this is automatic; in general posets
it is a genuine condition.
\end{remark*}

% ---------------------------------------------------------
% def:consistency
% ---------------------------------------------------------

\begin{definitionbox}{Definition}
\begin{definition}[Consistency]
\label{def:consistency}
Let $(S, \leq)$ be a partially ordered set.  Elements $x, y \in S$ are
\emph{consistent} (or \emph{compatible}) if they admit a common upper
bound:
\[
\exists z \in S, \quad x \leq z \;\land\; y \leq z.
\]
\end{definition}
\end{definitionbox}

\begin{remark*}[Standard quantified statement]
\[
\operatorname{Consistent}(x, y, S, \leq)
\;\coloneqq\;
\exists z \in S,\; x \leq z \land y \leq z.
\]
\end{remark*}

\begin{remark*}[Negated quantified statement]
\[
\neg\operatorname{Consistent}(x, y, S, \leq)
\;\coloneqq\;
\forall z \in S,\; x \not\leq z \lor y \not\leq z.
\]
\end{remark*}

\begin{remark*}[Interpretation]
Consistency asserts that $x$ and $y$ are compatible --- there is some
element of $S$ above both.  In a lattice where binary joins always exist,
consistency is automatic: take $z = x \vee y$.  In a general poset the
condition is nontrivial and may fail.  In domain theory, consistency
captures the idea that two partial computations describe information that
can be combined without contradiction.
\end{remark*}

\begin{dependencies}
\hyperref[def:lattice-order-partial-order]{Partial order}.
\end{dependencies}
